%------------------------------------------------------------------------------%
\documentclass[fleqn,utf8,aspectratio=169,12pt]{beamer}

\usepackage{gaal}

\title{Apresentação da Disciplina}

%------------------------------------------------------------------------------%
\begin{document}

\addtitle

%------------------------------------------------------------------------------%
\section{Introdução}

%------------------------------------------------------------------------------%
\begin{frame}{Disciplina}
   Disciplina: \emph{\large Geometia Analítica e Álgebra Linear}\\[7mm]
   Professor: \emph{\large Luis A. D'Afonseca}
\end{frame}

%------------------------------------------------------------------------------%
\begin{frame}{Disciplina}
\begin{itemize}[<+->]
  \item Aula semi invertida
  \item Estude o conteúdo com antecedência
  \item As apresentações são um resumo da teoria -- \emph{Estude o livro}
  \item Traga papel e lápis para fazer os exercícios em aula
  \item Estude o conteúdo seriamente entre uma aula e outra
  \item Não deixe o conteúdo acumular
  \item Não maratone o estudo de Matemática
  \item Preencha as lacunas do seu conhecimento de Matemática \\
        \href{https://pt.khanacademy.org}{Khan Academy}
\end{itemize}
\end{frame}

%------------------------------------------------------------------------------%
\section{Conteúdo}

%------------------------------------------------------------------------------%
\begin{frame}{Disciplina}
  \begin{enumerate}[<+->]
    \item Resolver sistemas lineares
    \item Realizar operações básicas envolvendo vetores
    \item Aplicar as técnicas vetoriais a problemas em geometria plana e espacial
    \item Representar e identificar retas, planos, cônicas por equações
    \item Determinar interseções, distâncias e ângulos entre retas e planos
    \item Identificar $R^2$ e $R^3$ como espaços vetoriais e seus subespaços.
    \item Determinar base e dimensão de subespaços de $R^2$ e $R^3$
    \item Aplicar processo de Gram-Schmidt para encontrar bases ortogonais
          e ortonormais de subespaços de $R^2$ e $R^3$
    \item Calcular autovalores e autovetores de uma matriz $2\times 2$ e $3\times 3$
  \end{enumerate}
\end{frame}

%------------------------------------------------------------------------------%
\section{Material para Estudo}

%------------------------------------------------------------------------------%
\begin{frame}{Material para Estudo}
  \begin{itemize}[<+->]
    \item Apostila Reginaldo \\[2mm]
          \href{https://www.dropbox.com/s/aa71ogpk8xski1j/gaalt1.pdf}%
               {https://www.dropbox.com/s/aa71ogpk8xski1j/gaalt1.pdf}
    \item Material de apoio do Pedro e Jane
    \item Notas de aula e provas anteriores \\[2mm]
          \href{\materialurl}{\materialurl}
  \end{itemize}
\end{frame}

%------------------------------------------------------------------------------%
\section{Lista Mínima}

%------------------------------------------------------------------------------%
\begin{frame}{Lista Mínima}
  \begin{enumerate}
    \item Estudar \href{https://pt.khanacademy.org}{Khan Academy}
%    \begin{enumerate}
%      \item \href{https://pt.khanacademy.org/math/precalculus}{Pré-cálculo}
%    \end{enumerate}
  \end{enumerate}

  \vfill
  Atenção: A prova é baseada no livro, não nas apresentações
\end{frame}

%------------------------------------------------------------------------------%
\end{document}
%------------------------------------------------------------------------------%
